\documentclass[10pt,a4paper]{amsart}
\usepackage[margin=19mm]{geometry}
\usepackage{amsmath,amssymb,mathtools}
\usepackage[scr=rsfs,cal=euler]{mathalfa}
\usepackage{xfrac}
\usepackage[unicode,hidelinks,bookmarks=false]{hyperref}
\newcommand{\ie}{i.e.\ }
\newcommand{\cf}{cf.\ }
\newcommand{\resp}{respectively\ }
\newcommand{\Subsection}{\subsection}
\providecommand{\nobreakdash}{\nobreak\hbox{-}}
\setlength{\emergencystretch}{2em}
\multlinegap0pt
\usepackage{bm}
\usepackage{bbm}
\usepackage{dsfont}
\usepackage{xspace}
\usepackage{cancel}
\usepackage{mathtools}
\usepackage{amsthm}
\makeatletter
\@ifundefined{theorem}{\newtheorem{theorem}{Theorem}}{}
\@ifundefined{lemma}{\newtheorem{lemma}[theorem]{Lemma}}{}
\makeatother
\newcommand{\supp}{\operatorname{supp}}
\title[Well-posedness of the Langmuir film problem]{Well-posedness of the Langmuir film problem}
\author{Yoichiro Mori, Shinya Okabe, Koya Sakakibara}
\date{Source: arXiv:2601.16482v1 (CC BY 4.0)}
\begin{document}
\maketitle

\noindent\emph{Source and licence.} Well-posedness of the Langmuir film problem. Version arXiv:2601.16482v1 (CC BY 4.0), distributed under the Creative Commons Attribution 4.0 International licence, as recorded in the arXiv licence metadata for this exact version and shown on the abstract page https://arxiv.org/abs/2601.16482v1. Full rights notice: licenses/arxiv-2601.16482-CC-BY-4.0.txt.\par\smallskip

\noindent\textit{Editorial scope.} Counted unit: the Lemma whose source label is \texttt{lem:resolvent\_estimate} in the article (source0.tex lines 1414--1423, source label \texttt{lem:resolvent\_estimate}), delivered together with the article's own complete original proof of it (source0.tex lines 1425--1596, one \texttt{proof} environment of 172 lines). No mathematics is written, repaired or re-derived: statement and proof are the article's own text line for line. Required context retained uncounted: none. Every definition, display and label of those lines is retained unchanged. Omitted from this package and not redistributed: 1--1413, 1424--1424, 1597--2413. Nothing inside a retained excerpt is omitted or altered: every retained body is delivered exactly as the article's lines are written, so no omission receipt of any kind accompanies this package. Citations: the retained excerpts carry no citation command at all, so no bibliography entry of the article is transcribed and no reference is redistributed. Reference adaptation: none was needed and none was made; every reference inside the delivered unit resolves inside it, so no unresolved reference or citation is produced. Printed slips kept literal: no printed word, symbol, hypothesis or number of the counted statement or of its proof was altered. Layout and class: the article's own class is \texttt{article} with packages including graphicx, mathtools, bm, amssymb, amsthm, amsmath, todonotes, geometry, cases, stmaryrd, xcolor, url; the wrapper is the lane's pinned \texttt{amsart} editorial wrapper at 10pt on A4, which restates the theorem-like environments the retained text needs and adds the article's own macro definitions that the retained text needs (\texttt{\textbackslash supp}), with extra wrapper packages (bm, bbm, dsfont, xspace, cancel, mathtools). The delivered numbering is the unit's own. Assets: no retained span contains a figure environment, a graphics-inclusion command, a PDF-page inclusion, a table or a TikZ picture; no image file is copied into this package, the package declares no asset dependency and no image file is redistributed with it. Licence: version arXiv:2601.16482v1 of this article is distributed under the Creative Commons Attribution 4.0 International licence, as recorded in the arXiv licence metadata for this exact version and shown on the abstract page https://arxiv.org/abs/2601.16482v1. A full rights notice is delivered at licenses/arxiv-2601.16482-CC-BY-4.0.txt. Characters: every retained excerpt is pure ASCII, so the delivered engine, XeLaTeX with its default Latin Modern text font, reports no missing character.
\begin{lemma}
    \label{lem:resolvent_estimate}
    Let $\eta\in V$.
    Then, there exist a constant $\lambda\ge0$ and a positive constant $C_{\delta,\eta}$, depending only on $\delta$ and $\eta$, such that
    \begin{align}
        \|(z+\lambda-A_s(\eta))\gamma\|_{H^1} \ge C_{\delta,\eta}\left(|z|\|\gamma\|_{H^1} + \|\gamma\|_{H^3}\right)
        \label{eq:resolvent_estimate}
    \end{align}
    holds for all $z\in\Sigma_{0,\delta}$, $\gamma\in H^3$, and $s\in[0,1]$.
\end{lemma}

\begin{proof}
We divide the proof into three steps.

\textit{Step 1}.
We first introduce a partition of unity on $\mathbb{S}^1$.
Define a family $\{\phi_k\}_{k=1}^N$ of functions on $\mathbb{S}^1$ by
\begin{align*}
    \phi_k(\cdot)\coloneqq\phi(\cdot-x_k),\qquad
    x_k=\frac{k-1}{N},
\end{align*}
where $\phi\in C^\infty(\mathbb{S}^1)$ satisfies
\begin{align}
    \phi\ge0\quad\text{on }\mathbb{S}^1,\qquad
    \supp\phi\subset\left[-\frac{1}{N},\frac{1}{N}\right],\qquad
    \sum_{k=1}^N\phi_k^2\equiv1\quad\text{on }\mathbb{S}^1.
    \label{eq:partition-of-unity}
\end{align}
It is straightforward to verify the identity
\begin{align}
    \sum_{k=1}^N\|\phi_k u\|_{L^2}^2 = \|u\|_{L^2}^2
    \label{eq:sum_phi_k-u_L2}
\end{align}
for all $u\in L^2$.
Furthermore, there exists a positive constant $C_{\phi,N}$, depending only on $\phi$ and $N$, such that the following estimates hold:
\begin{align}
    \sum_{k=1}^N\|\phi_k'\|_{L^\infty}^2\le C_{\phi,N},\qquad
    \sum_{k=1}^N\|\phi_k''\|_{L^\infty}^2 \le C_{\phi,N}.
    \label{eq:phi'_phi''_estimate}
\end{align}

\textit{Step 2}.
In this step, we show that for any $\varepsilon>0$, there exist positive constants $C_{\varepsilon,2}$ and $C_{\varepsilon,\infty}$, depending only on $\varepsilon$, such that the estimates
\begin{align}
    \|v'\|_{L^2}^2 \le C_{\varepsilon,2}\|v\|_{L^2}^2 + \varepsilon\|v''\|_{L^2}^2,\qquad
    \|v'\|_{L^\infty}^2 \le C_{\varepsilon,\infty}\|v\|_{L^2}^2 + \varepsilon\|v''\|_{L^2}^2
    \label{eq:v'_L2_Linf}
\end{align}
hold for all $v\in H^2$.
In particular, we can choose $C_{\varepsilon,2}=1/(4\varepsilon)$ and $C_{\varepsilon,\infty}=27/(16\varepsilon^3)$.

Integrating by parts and applying the Cauchy--Schwarz and Young inequalities, we obtain
\begin{align}
    \|v'\|_{L^2}^2
    &=-\int_{\mathbb{S}^1}v\cdot v''\,\mathrm{d}x
    \le\|v\|_{L^2}\|v''\|_{L^2}
    \le\frac{1}{4\varepsilon}\|v\|_{L^2}^2 + \varepsilon\|v''\|_{L^2}^2.
    \label{eq:first_ineq}
\end{align}
Thus, setting $C_{\varepsilon,2}\coloneqq1/(4\varepsilon)$, we obtain the first desired inequality.

Next, consider the second inequality.
Since $\int_{\mathbb{S}^1}v'=0$, the one-dimensional Gagliardo--Nirenberg interpolation inequality, together with \eqref{eq:first_ineq} and Young's inequality, yields
\begin{align*}
    \|v'\|_{L^\infty}^2
    \le2\|v'\|_{L^2}\|v''\|_{L^2}
    \le2\left(\|v\|_{L^2}\|v''\|_{L^2}\right)^{1/2}\|v''\|_{L^2}
    =2\|v\|_{L^2}^{1/2}\|v''\|_{L^2}^{3/2}
    \le\frac{27}{16\varepsilon^3}\|v\|_{L^2}^2 + \varepsilon\|v''\|_{L^2}^2,
\end{align*}
which is precisely the second inequality with $C_{\varepsilon,\infty}\coloneqq27/(16\varepsilon^3)$. This completes Step~2.

Note that there exists a positive constant $C_\delta$, depending only on $\delta$, such that
\begin{align}
    |z+a|\ge C_\delta(|z|+a)
    \label{eq:abs_below}
\end{align}
holds for all $z\in\Sigma_{0,\delta}$ and $a\ge0$.
In particular, we can choose $C_\delta=\sqrt{(1-\cos\delta)/2}$.

For $z\in\Sigma_{0,\delta}$, $\lambda\ge0$, $s\in[0,1]$, and $\gamma\in H^3$, set
\begin{align*}
    f\coloneqq(z+\lambda)\gamma-A_s(\eta)\gamma=(z+\lambda)\gamma - B_s(\eta)\gamma''.
\end{align*}
Differentiating this equation, we obtain
\begin{align}
    f'=(z+\lambda)v - B_s(\eta)'v' - B_s(\eta)v'',
    \label{eq:f'}
\end{align}
where $v\coloneqq\gamma'$.

\textit{Step 3}.
In this step, we show that there exists a constant $C_{\delta,\eta}$, depending only on $\delta$ and $\eta$, such that
\begin{align*}
    \|f\|_{H^1}^2 \ge C_{\delta,\eta}^2\left(|z|\|\gamma\|_{H^1} + \|\gamma\|_{H^3}\right)^2
\end{align*}
for sufficiently large $\lambda$.
Let $B_k\coloneqq B(\eta)|_{x=x_k}$.
Multiplying both sides of \eqref{eq:f'} by $\phi_k$ gives
\begin{align*}
    (z+\lambda)\phi_kv - B_k(\phi_kv)''
    =\phi_kf'+B_s(\eta)'\phi_kv'+\phi_k(B_s(\eta)-B_k)v''-2B_k\phi_k'v'-B_k\phi_k''v
    \eqqcolon b_k.
\end{align*}
Applying \eqref{eq:abs_below} and the Cauchy--Schwarz inequality, we obtain
\begin{align*}
    \|b_k\|_{L^2}^2
    &= \|(z+\lambda)\phi_kv - B_s(\eta)(\phi_kv)''\|_{L^2}^2
    =\sum_{n\in\mathbb{Z}}|z+\lambda+B_k(2\pi n)^2|^2|(\phi_kv)_n|^2\\
    &\ge\sum_{n\in\mathbb{Z}}C_\delta^2(|z|+\lambda+B_k(2\pi n)^2)^2|(\phi_kv)_n|^2
    \ge C_\delta^2\left[(|z|+\lambda)^2\|\phi_kv\|_{L^2}^2 + B_\ast^2\|(\phi_kv)''\|_{L^2}^2\right].
\end{align*}
Using the identity $(\phi_kv)'' = \phi_kv''+2\phi_k'v'+\phi_k''v$, Young's inequality implies
\begin{align*}
    \|(\phi_kv)''\|_{L^2}^2
    \ge\frac{1}{2}\|\phi_kv''\|_{L^2}^2 - 13\|\phi_k'v'\|_{L^2}^2 - 7\|\phi_k''v\|_{L^2}^2.
\end{align*}
The estimates \eqref{eq:phi'_phi''_estimate} and \eqref{eq:v'_L2_Linf} imply that
\begin{align*}
    \sum_{k=1}^N\|\phi_k'v'\|_{L^2}^2
    \le C_{\phi,N}\left(C_{\varepsilon,2}\|v\|_{L^2}^2 + \varepsilon\|v''\|_{L^2}^2\right),\qquad
    \sum_{k=1}^N\|\phi_k''v\|_{L^2}^2 \le C_{\phi,N}\|v\|_{L^2}^2.
\end{align*}
Combining these with \eqref{eq:sum_phi_k-u_L2}, we have
\begin{align}
    \sum_{k=1}^N\|b_k\|_{L^2}^2
    &\ge
    C_\delta^2\left[(|z|+\lambda)^2 - 13B_\ast^2C_{\phi,N}C_{\varepsilon,2} - 7B_\ast^2C_{\phi,N}\right]\|v\|_{L^2}^2 + C_\delta^2B_\ast^2\left(\frac{1}{2}-13C_{\phi,N}\varepsilon\right)\|v''\|_{L^2}^2.
    \label{eq:sum_bk_lower}
\end{align}
Next, we consider an upper bound for $\sum_{k=1}^N\|b_k\|_{L^2}^2$.
It follows from \eqref{eq:sum_phi_k-u_L2} and \eqref{eq:v'_L2_Linf} that
\begin{align*}
    \sum_{k=1}^N\|\phi_k(f'+B_s(\eta)'v')\|_{L^2}^2
    &=\|f'+B_s(\eta)'v'\|_{L^2}^2
    \le2\|f'\|_{L^2}^2 + 2\|B_s(\eta)'\|_{L^2}^2\|v'\|_{L^\infty}^2\\
    &\le2\|f'\|_{L^2}^2 + 2\|B(\eta)'\|_{L^2}^2\left(C_{\varepsilon,\infty}\|v\|_{L^2}^2 + \varepsilon\|v''\|_{L^2}^2\right).
\end{align*}
Since $\phi_k$ is supported in $[x_k-N^{-1},x_k+N^{-1}]$, it follows from \eqref{eq:sum_phi_k-u_L2} that
\begin{align*}
    \sum_{k=1}^N\|\phi_k(B_s(\eta)-B_k)v''\|_{L^2}^2
    &\le\sum_{k=1}^N\sup_{|x-x_k|\le N^{-1}}|B_s(\eta)-B_k|^2\|\phi_kv''\|_{L^2}^2
    \le\omega_{B,N}\|v''\|_{L^2}^2,
\end{align*}
where we set $\omega_{B,N}\coloneqq\sup_{|x-y|\le N^{-1}}|B(\eta(x))-B(\eta(y))|^2$.
Therefore, using Young's inequality and \eqref{eq:phi'_phi''_estimate}, we have
\begin{align*}
    \sum_{k=1}^N\|b_k\|_{L^2}^2
    &\le\sum_{k=1}^N\left[4\|\phi_k(f'+B_s(\eta)'v')\|_{L^2}^2 + 4\|\phi_k(B_s(\eta)-B_k)v''\|_{L^2}^2 + 16\|B_k\phi_k'v'\|_{L^2}^2 + 4\|B_k\phi_k''v\|_{L^2}^2\right]\\
    &\le8\|f'\|_{L^2}^2 + \left[8C_{\varepsilon,\infty}\|B(\eta)'\|_{L^2}^2 + 4C_{\phi,N}\|B(\eta)\|_{L^\infty}^2\left(4C_{\varepsilon,2} + 1\right)\right]\|v\|_{L^2}^2\\
    &\quad+\left[4\omega_{B,N} + \left(8\|B(\eta)'\|_{L^2} + 16C_{\phi,N}\|B(\eta)\|_{L^\infty}^2\right)\varepsilon\right]\|v''\|_{L^2}^2.
\end{align*}
Hence, combining with \eqref{eq:sum_bk_lower}, we obtain
\begin{align*}
    8\|f'\|_{L^2}^2
    &\ge(|z|+\lambda)^2\left(C_\delta^2 - \frac{B_\ast^2C_{\phi,N}C_\delta^2(13C_{\varepsilon,2}+7) + 8C_{\varepsilon,\infty}\|B(\eta)'\|_{L^2}^2 + 4C_{\phi,N}\|B(\eta)\|_{L^\infty}^2(4C_{\varepsilon,2}+1)}{\lambda^2}\right)\|v\|_{L^2}^2\\
    &\quad+\left(\frac{1}{2}C_\delta^2B_\ast^2 - 4\omega_{B,N} - \left(13B_\ast^2C_{\phi,N}C_\delta^2 + 8\|B(\eta)'\|_{L^2}^2 + 16C_{\phi,N}\|B(\eta)\|_{L^\infty}^2\right)\varepsilon\right)\|v''\|_{L^2}^2.
\end{align*}
Since $B(\eta)\in H^1\hookrightarrow C^\alpha$ with $\alpha\in[0,1/2)$ by Sobolev embedding, $B(\eta)$ is, in particular, uniformly continuous.
Therefore, taking $N$ sufficiently large and fixing it, we can ensure that $4\omega_{B,N} \le C_\delta^2B_\ast^2 / 8$.
Then, by choosing $\varepsilon$ sufficiently small and fixing it, we can ensure that the coefficient of $\|v''\|_{L^2}^2$ is bounded from below by $\frac{1}{4}C_\delta^2B_\ast^2$.
Accordingly, by choosing $\lambda\ge1$ sufficiently large, we see that the coefficient of $\|v\|_{L^2}^2$ is bounded from below by $C_\delta^2/2$.
Thus, there exists a constant $C_{\delta,\eta}$, depending only on $\delta$ and $\eta$, such that
\begin{align}
    \|f'\|_{L^2}^2 \ge 2C_{\delta,\eta}^2\left((|z|+\lambda)^2\|v\|_{L^2}^2 + \|v''\|_{L^2}^2\right)
    \label{eq:f'_lower}
\end{align}
holds.
In a similar but simpler manner (working directly with $f$ instead of $f'$), we can also show that
\begin{align}
    \|f\|_{L^2}^2 \ge 2C_{\delta,\eta}^2\left((|z|+\lambda)^2\|\gamma\|_{L^2}^2 + \|\gamma''\|_{L^2}^2\right)
    \label{eq:f_lower}
\end{align}
holds by choosing $C_{\delta,\eta}$ smaller if necessary.
Noting that $v=\gamma'$ and thus $\|v''\|_{L^2} = \|\gamma'''\|_{L^2}$, we combine \eqref{eq:f'_lower} with \eqref{eq:f_lower} to obtain
\begin{align*}
    \|f\|_{H^1}^2
    &\ge2C_{\delta,\eta}^2\left[(|z|+\lambda)^2(\|\gamma\|_{L^2}^2 + \|\gamma'\|_{L^2}^2) + \|\gamma''\|_{L^2}^2 + \|\gamma'''\|_{L^2}^2\right]\\
    &\ge2C_{\delta,\eta}^2\left(|z|^2\|\gamma\|_{H^1}^2 + \|\gamma\|_{H^3}^2\right)
    \ge C_{\delta,\eta}^2\left(|z|\|\gamma\|_{H^1} + \|\gamma\|_{H^3}\right)^2,
\end{align*}
where we have used $\lambda\ge1$ in the second inequality.
\end{proof}

\end{document}
